%%%%%%%%%%  scan idx 685  |  folio 672  %%%%%%%%%%
\SGAPrintPage{672}%
SGA 4 XVI 3.7), and, conjecturally, when $Y$ and $X$ are both regular (by the absolute cohomological purity theorem, proved at least for excellent $X$ of characteristic zero (SGA 4 XIX 3.2)). When $X$ and $Y$ are smooth over the field of complex numbers, one can also use the integral Chern class on the locally compact space $X(\C)$ of complex points of $X$. Formula (3.1) is established in this case by \uline{Atiyah} and \uline{Hirzebruch} [1], who work more generally with an immersion of complex analytic varieties (using the ``naive'' $K^\bullet$); it follows that the $\ell$-adic formula (3.1) holds whenever $Y$ and $X$ are smooth over an \uline{algebraically closed field of characteristic zero}, using the ``Lefschetz principle'' and the general results of SGA 4. This therefore makes the validity of the cohomological formula (3.1) extremely plausible; but observe that it has not been proved at present even when $X$ and $Y$ are smooth over a field $k$ of characteristic zero, when $k$ is not algebraically closed.

\sgaStableTarget{sga6.e14.para.3-4}\textbf{3.4.} One can also work with the cohomological theory of Chern classes in the De Rham or Hodge style [13]. The question of the validity of (3.1) in this setting arises only when $X$ is not of characteristic zero, since in any case (no.~6) the formula in question is known to hold modulo torsion. For smooth quasi-projective schemes over a field $k$, the De Rham or Hodge version of formula (3.1) would moreover follow from the cycle-class version of formula (3.1) mod. algebraic equivalence, so that the De Rham or Hodge cohomological version gives another accessible way to test the validity of the ``cycle-class'' version of (3.1).

\sgaStableTarget{sga6.e14.para.3-5}\textbf{3.5.} The proof [1] of the transcendental cohomological version of (3.1) is essentially transcendental in nature, using tubular neighbourhoods of $Y$ in $X$. This suggests that a proof of a formula of type (3.1) might be connected with the systematic consideration of groups $K^\bullet$ and cohomology groups \uline{with support in} $Y$, which might provide a sufficient technical means of ``localization around $Y$''.

\sgaStableTarget{sga6.e14.sec.4}\section*{4. Relations between $K^\bullet(X)$ and the Chow ring $A(X)$}

\sgaStableTarget{sga6.e14.para.4-1}\textbf{4.1.} For simplicity, suppose that $X$ is a smooth quasi-projective scheme over a field $k$ (cf no.~8 for a more general case), so that we
%%%%%%%%%%  scan idx 686  |  folio 673  %%%%%%%%%%
\SGAPrintPage{673}%
have the Chow ring $A(X)$ of cycle classes modulo rational equivalence [4]. As variants of this ring, which is moreover defined under substantially more general considerations, there are the rings $\operatorname{Gr}^\bullet(X)$ (the associated graded of the $\lambda$-filtration of $K^\bullet(X)$) and $\Gtop^\bullet(X)$ (the associated graded of the filtration of $K^\bullet(X)\simeq K_\bullet(X)$ by the codimension of supports). There are theories of Chern classes with values in each of these three rings (cf. [9] for $A(X)$, and this seminar for $\operatorname{Gr}^\bullet(X)$, which maps to $\Gtop^\bullet(X)$). There is also a cohomological theory ($\ell$-adic, say, with $\ell$ prime to the characteristic of $k$) of Chern classes (SGA 5 VII):
\[
 c_i(\underline{E})\in H^{2i}(X,\underline{Z}_\ell(i)),
\]
(where $\underline{E}$ is a vector bundle on $X$, or indeed any element of $K^\bullet(X)$). These cohomological theories are related by the canonical homomorphisms (compatible with the multiplicative structures)
\sgaStableTarget{sga6.e14.eq.4-1}\begin{equation}\tag{4.1}
\begin{tikzcd}[column sep=large,row sep=large]
 H^{2*}(X,\underline{Z}_\ell(*)) & {} \\
 A(X) \arrow[u,"\gamma"] \arrow[r,"\phi"] & \Gtop^\bullet(X) \\
 {} & \operatorname{Gr}^\bullet(X), \arrow[u,"j"]
\end{tikzcd}
\end{equation}
where we have set
\[
 H^{2*}(X,\underline{Z}_\ell(*))=\coprod_i H^{2i}(X,\underline{Z}_\ell(i)).
\]
The homomorphism $j$ comes from the fact that the $\lambda$-filtration of $K^\bullet(X)$ is finer than the topological filtration by the codimension of the support (Exp. X), and it is an isomorphism modulo torsion (Exp. VII 4.11). The homomorphism $\phi$ is defined, for example, in [RRR], and associates to the class of the cycle of codimension $i$ associated with a closed subscheme $Y$ of $X$ of codimension $i$ the element $\mathrm{c}\ell^\bullet(\underline{O}_Y)$ of $\operatorname{Filt}^i_{\mathrm{top}}(K^\bullet(X))$ mod $\operatorname{Filt}^{i+1}_{\mathrm{top}}(K^\bullet(X))$. It is a surjective homomorphism, which we shall also see (4.2) to be an isomorphism modulo torsion. Thus, \uline{modulo torsion the three rings} $A(X)$, $\Gtop^\bullet(X)$ \uline{and} $\operatorname{Gr}^\bullet(X)$ \uline{are canonically isomorphic}, so \uline{a posteriori} the Riemann--Roch formulas expressed using one or another of these rings (tensored with $\Q$) are equivalent. The graded ring $\operatorname{Gr}^\bullet(X)$, used particularly in this seminar as a substitute for the Chow ring, has the appreciable advantage of being defined for every locally ringed topos and, in particular, for every scheme
%%%%%%%%%%  scan idx 687  |  folio 674  %%%%%%%%%%
\SGAPrintPage{674}%
(regular or otherwise). It has, on the other hand, the disadvantage (4.6) relative to $A(X)$, $\Gtop^\bullet(X)$ that, as it stands, it is not covariant with respect to proper morphisms (even for smooth quasi-projective schemes over a field); it acquires this property only after tensoring with $\Q$. Another advantage of $A(X)$ over its competitors $\Gtop^\bullet$ and $\operatorname{Gr}^\bullet(X)$ is the existence of a homomorphism from $A(X)$ to $H^{2*}(X,\underline{Z}_\ell(*))$ that allows cohomological Chern classes to be interpreted as images of Chern classes with values in $A(X)$. In general this homomorphism does not factor through $\Gtop^\bullet(X)$ (4.7), and it is at least doubtful that there is in general a homomorphism from $\operatorname{Gr}^\bullet(X)$ to $H^{2*}(X,\underline{Z}_\ell(*))$ taking Chern classes to Chern classes (Cf. no.~5). Of course, these difficulties disappear after tensoring with $\underline{Q}$, and one obtains a homomorphism
\sgaStableTarget{sga6.e14.eq.4-2}\begin{equation}\tag{4.2}
 \Gtop^\bullet(X)\longrightarrow H^{2*}(X,\underline{Q}_\ell(*))=
 H^{2*}(X,\underline{Z}_\ell(*))\otimes_{\underline{Z}_\ell}\underline{Q}_\ell
\end{equation}
taking Chern classes to Chern classes.

\sgaStableTarget{sga6.e14.para.4-2}\textbf{4.2.} To establish that the surjective homomorphism $\phi$ of (4.1) is an isomorphism modulo torsion, one uses the Chern homomorphism (deduced from the theory of Chern classes [9])
\[
 K^\bullet(X)\longrightarrow \Ch(A(X))
\]
(where the second member is the Chern ring of $A(X)$ defined in Exp.~V), which is, as is easily seen, compatible with the topological filtration of $K^\bullet(X)$ and the natural filtration of $\Ch(A(X))$, and gives, on passing to the associated graded groups, a homomorphism of graded groups
\sgaStableTarget{sga6.e14.eq.4-3}\begin{equation}\tag{4.3}
 \psi:\Gtop^\bullet(X)\longrightarrow A(X).
\end{equation}
With this understood, we shall prove the formulas
\sgaStableTarget{sga6.e14.eq.4-4}\begin{equation}\tag{4.4}
\left\{
\begin{aligned}
 \psi\phi(x)&=(-1)^{i-1}(i-1)!\,x\quad \text{mod. torsion}\quad \text{for } x\in A^i(X)\\
 \phi\psi(x)&=(-1)^{i-1}(i-1)!\,x\quad \text{mod. torsion}\quad \text{for } x\in \Gtop^i(X),
\end{aligned}
\right.
\end{equation}
the first of which moreover implies the second, since $\phi$ is surjective. These formulas clearly imply that $\phi$ and $\psi$ are isomorphisms modulo torsion and that, up to a factor $(-1)^{i-1}(i-1)!$ in degree $i$, they are inverses of one another.

To prove the first formula (4.4), one can clearly suppose that $x$ is defined by an irreducible closed subscheme $Y$ of $X$ of codimension $i$, in which case the formula is equivalent to
%%%%%%%%%%  scan idx 688  |  folio 675  %%%%%%%%%%
\SGAPrintPage{675}%
\sgaStableTarget{sga6.e14.eq.4-5}\begin{equation}\tag{4.5}
 c_i(\underline{O}_Y)=(-1)^{i-1}(i-1)!\,\cl(Y)\quad \bmod
\text{ torsion},
\end{equation}
where in the first member $\underline{O}_Y$ is regarded as a coherent sheaf on $X$, thus defining an element of $K^\bullet(X)$, and in the second member $\mathrm{c}\ell(Y)$ denotes the class of $Y$ in $A^i(X)$. Since $A^i(X)$ does not change on removing from $X$ a closed subset of codimension $>i$, we are reduced to the case in which $Y$ is regular, hence smooth over the base field $k$ when the latter is perfect. But in this latter case, formula (4.5) is precisely the Riemann--Roch formula for the inclusion $Y\to X$ and the element $1$ of $K^\bullet(Y)$, written using the theory of Chern classes with values in the Chow ring $A(X)$. This formula is established in [2] under the conditions considered here, by a method very close to the one used in this seminar for the analogous formula using $\operatorname{Gr}^\bullet(X)$. This therefore establishes, at least for perfect $k$, formula (4.4) and consequently the fact that $\phi$ is an isomorphism mod. torsion (which in turn implies, a posteriori, that the Riemann--Roch formula ``with denominators'' for a proper morphism of smooth algebraic schemes over $k$ is essentially the same whether written using $A(X)$ as in [2] or using $\operatorname{Gr}^\bullet(X)$ as in this seminar!). The case with no hypothesis on $k$ reduces easily to the case in which $k$ is algebraically closed, using the fact that, for a finite extension $k'$ of degree $n$ of $k$, the kernel of the natural map $A(X)\xrightarrow{f^*}A(X\otimes_k k')$ is annihilated by $n$ (as is seen by using the trace homomorphism $f_*:A(X\otimes_k k')\to A(X)$, satisfying $f_*f^*=n\,\id$).\SGAAuthorFootText{*}{It is therefore established, thanks to P. JOUANOLOU, for $X$ smooth and quasi-projective over a field (cf. note p.5).}

\sgaStableTarget{sga6.e14.para.4-3}\textbf{4.3.} It seems quite plausible that formulas (4.4) hold as they stand and not only modulo torsion, or, equivalently, that formula 4.5. holds not only modulo torsion. This question is a very special case of the validity in $A(X)$ of the ``Riemann--Roch formula without denominators'' (3.1) for an immersion $Y\to X$ (namely, the case $i=d$). The validity of this special case may be regarded as quite plausible. Indeed, we already pointed out in no.~3 that, when $k$ is of characteristic zero, the Riemann--Roch formula without denominators holds at least in $\Gr^{\bullet}_{\mathrm{top}}(X)$; this implies in that case that both members of (4.5) have the same image in $\Gr^{\bullet}_{\mathrm{top}}(X)$, or equivalently that the second of formulas (4.4) holds; this result will extend of itself to arbitrary characteristic once
%%%%%%%%%%  scan idx 689  |  folio 676  %%%%%%%%%%
\SGAPrintPage{676}%
the question raised in no.~1 has been satisfactorily resolved. In support of formula (4.5) in $A(X)$ itself, observe that it is clear a priori that the restriction of the first member to $X-Y$ is zero; thus, with $Y$ assumed irreducible, the first member must a priori be an integer multiple of $\mathrm{c}\ell(Y)$:
\[
 c_i(Y)=n(Y)\mathrm{c}\ell(Y).
\]
This already proves the exact formula (4.5) when $\mathrm{c}\ell(Y)$ is not a torsion element of $A^i(X)$. It is immediate that this is so when $X$ is projective (and not merely quasi-projective) over the field $k$: note that for a polarization $\xi\in A^1(X)$ of $X$, one has $\deg(Y\cdot\xi^{\dim Y})>0$; this proves the desired formula in this case. Consequently, this formula also holds whenever $X$ is isomorphic to an open subscheme of a smooth projective scheme over $k$, which will be the case when a suitable solution to the resolution-of-singularities problem is available. This is so, in particular, when $k$ has characteristic zero, by Hironaka [11].

The validity of the exact formula (4.5), hence of (4.4), would imply that the kernel of $\phi^i:A^i(X)\longrightarrow \Gtop^i(X)$ is annihilated by $(i-1)!$, and in particular \uline{that} $\phi^2:A^2(X)\to\Gtop^2(X)$ \uline{is an isomorphism}. (That $\phi^1$ is an isomorphism already follows immediately from X 5.3.2.) Unless I am mistaken, the fact that $\phi^2$ is an isomorphism has not been proved at present, even in the special case $\dim X=3$. In any event, the writer was hasty in asserting the validity of (4.4) and (4.5) in [9, no.~4, 3\textsuperscript{o}].

\sgaStableTarget{sga6.e14.para.4-4}\textbf{4.4.} At present, the main advantage of $\Gtop^\bullet(X)$ (or $\operatorname{Gr}^\bullet(X)$) over the Chow ring $A(X)$ is that it is better suited to explicit computation of certain intersection formulas between cycles, in the case of ``excess'' intersections, i.e. of excessively large dimension. Indeed, defining the intersection class by means of Chow's lemma, which allows a cycle to be moved sufficiently within its class, does not readily lend itself to computation. A striking example is the following fundamental formula [Exp. VII 2.7] for a closed immersion $f:Y\to X$ of smooth quasi-projective schemes over $k$:
\sgaStableTarget{sga6.e14.eq.4-6}\begin{equation}\tag{4.6}
 f^*f_*(y)=y\,c_d(\underline{N}),
\end{equation}
a formula valid in $\Gtop^\bullet(Y)$, with $y$ in $\Gtop^\bullet(Y)$ and $\underline{N}$ the conormal sheaf
